\newcommand{\inputtag}[1]{**\textless #1\textgreater**}



\section{Dataset Construction}\label{ch2:sec:the-dataset}

% [CH2-010][Original Start]
% In this section, we demonstrate the process of constructing the training dataset. The raw textual data are derived from the Federal Open Market Committee (FOMC) meeting minutes, complemented by corresponding macroeconomic and financial market indicators. These data are then systematically processed and transformed into structured training materials for model development.
% [CH2-010][Original End]
This section describes the construction of the training corpus. The raw text comes from Federal Open Market Committee (FOMC) Minutes and is paired with macroeconomic and financial indicators. We then process these sources into structured samples for model training and evaluation.


\subsection*{FOMC Minutes and Indicators}\label{ch2:sec:raw_data}
% Theseminutes were written from a third-person perspective and often included many specific details,attributing specific policy views to individual FOMC members

% [CH2-011][Original Start]
% The Federal Open Market Committee (FOMC) meetings have been recorded since the committee's formation under the Banking Act of 1935 (\cite{todd2016corollary}). Initially, from 1936 to 1967, the minutes were released only once a year in these early years. Between 1967 and 1992, the minutes were named as ``the Minutes of Actions''. combined with ``the Record of Policy Actions'', these documents provided detailed insights not only into the policies but also the reasoning and proceedings behind them. Starting in 1993, a timely record of all policy issues and decisions made by the Committee was published after each meeting, although the documentation lacked a structured format. After 2009, the minutes were organized into distinct sections, including "Economic Situation," "Financial Situation," "Economic Outlook," and "Committee Policy Action," and the Committee has continued to use this structure to the present day. Therefore, we have chosen to use the minutes post-2009 as our primary source material.
% [CH2-011][Original End]
FOMC meetings have been documented since the Committee's establishment under the Banking Act of 1935 (\cite{todd2016corollary}). Publication practices evolved over time: early records were less frequent and less standardized, while post-2009 Minutes adopted a relatively stable section structure. Because section-level consistency is essential for supervised alignment, we use post-2009 Minutes as the primary source corpus.



% First of all, we downloaded the data from 28th Jan 2009 to 29th Jan 2025 from the Federal Reserve System Website and named it as our \textbf{main dataset}. For FOMC Minutes, there are 128 documents in total. These Minutes are well tilted with 5 sub-sections, ``Staff Review of the Economic Situation'', ``Staff Review of the Financial Situation'', ``Staff Economic Outlook'', ``Participants' Views on Current Conditions and Economic Outlook'', and ``Committee Policy Actions''. There are 486 sections in total. Although some other sections may be included in the Minutes, for instance, there is a ``Developments in Financial Markets and the Federal Reserve`s Balance Sheet'' in the 20090624 Minutes, these five sections are all included in the Minutes after 2009. We also downloaded the Minutes before 2009 (3rd Feb 1993 - 12nd Dec 2008), but these documents do not contain a clear title for different sections. Hence, we named it as our \textbf{supplementary dataset} and conducted a further process on these Minutes in next section .
%
% On average, each Minutes contains 15,119 tokens. There are 1,401,485 tokens in the dataset. Tokens are integer numbers representing different ``words'', which serve as inputs for the model. 
% [Codex 修改版] 2026-02-26: 统一口径并避免与 Table~\\ref{tab:ch2:summary_statis} 的统计不一致；用“语料+分段样本”描述替代相互矛盾的总量数字。
% [CH2-012][Original Start]
% First, we downloaded FOMC Minutes spanning January 2009 to January 2025 from the Federal Reserve website and treated them as our \textbf{main text corpus}. The post-2009 Minutes follow a largely stable section structure; we therefore parsed the major sections (e.g., staff reviews, staff outlook, participants' views, and policy actions) when available, yielding 486 section-level samples in total. We also downloaded Minutes prior to 2009 (February 1993 to December 2008). Because these earlier documents lack consistent section titles, we treat them as a \textbf{supplementary corpus} and apply additional processing steps described below.
% [CH2-012][Original End]
% [Codex 修改版] 2026-04-17: 用当前 raw-data 审计补充会议总数口径，同时保留后面的旧统计表作为历史快照。
We downloaded Minutes from January 2009 to January 2025 from the Federal Reserve website as the \textbf{main text corpus}. Given the stable section structure after 2009, we parsed major sections (staff reviews, staff outlook, participants' views, and policy actions) and obtained 486 section-level samples. A current raw-data audit of \texttt{dataset/raw\_data/prompt\_summary\_20250512.jsonl} recovers 128 unique meeting dates over this span. We also collected Minutes from February 1993 to December 2008 as a \textbf{supplementary corpus}. Because these earlier files lack stable section headings, they require additional processing described below.

Table~\ref{tab:ch2:summary_statis} reports token-level summary statistics for both full Minutes and extracted sections under the tokenizer used in our experiments. Across the sample, Minutes length ranges from roughly 8,000 to 25,000 tokens, while most individual analytical sections fall below 4,096 tokens.

% [CH2-013][Original Start]
% Notably, the term ``word'' in this context does not strictly refer to whole words. Instead, LLMs operate on ``sub-words'' or ``tokens'', which are components of words. For instance, ``Negative'' might be tokenized as ``Neg'' and ``agtive'' to reduce the size of the dictionary, which is known as token dictionary or tokenizer. Each LLM has its unique tokenizer. More powerful LLMs typically have larger tokenizers. For instance, ChatGPT-4's tokenizer contains 256,000 tokens, while LLAMA3's tokenizer includes 128,000 tokens. Table \ref{tab:ch2:summary_statis} provides the summary statistics for the tokens.
% [CH2-013][Original End]
In this chapter, ``word'' does not necessarily mean a full lexical word. LLMs operate on tokens (often sub-word units). For example, ``Negative'' may be split into smaller units such as ``Neg'' and ``ative''. This tokenization is controlled by the model-specific tokenizer. Larger models often use larger vocabularies; for example, GPT-4 class models use larger token vocabularies, while LLaMA 3 uses a 128,000-token vocabulary. Table~\ref{tab:ch2:summary_statis} reports token-level summary statistics.

% Notably, LLM utilize sub-words rather than entire words. Sub-words are segments of words, such as ``tive'', ``sion'', and ``ing'', which help increase the model's efficiency. The transaction from a ``word'' to some ``tokens'' is facilitated by the tokenizer, which functions similarly to a dictionary of tokens and sub-words. Generally, a more powerful model will have a larger tokenizer. For instance, the LLAMA2's tokenizer has 8K entries, while LLAMA3's contains 128K. 

\begin{table}[htbp]
    \scriptsize
    \centering
    \begin{tabular}{p{3.8cm}cccccccc}
        \toprule
        Names&                                                Count&	    Mean&	    Std&	 Min&	   25\%&	  50\%&	   75\%&	Max \\
        \hline
        Total Minutes&	                                        123&	15119.50&	5136.89&	8144&	  10545&	 13117&	  20234&	25106 \\
        Total Sections&	                                        486&	 2883.71&	2921.00&	 294&	 1391.5&	1729.5&	3014.75&	13621\\
        Staff Economic Outlook&	                                123&	 3002.20&	 779.20&	 294&	   2556&	  3076&	   3481&	5080 \\
        Staff Review of the Economic Situation&	                120&	 1348.81&	 314.13&	 653&	1178.25&	  1356&	   1490&	2929 \\
        Staff Review of the Financial Situation&	            119&	 1491.61&	 291.48&	 903&	   1293&	  1479&	 1643.5&	2461 \\
        Committee Policy Action&	                            122&	 5625.28&	4646.81&	1262&	   1770&	2448.5&	  11118&	13621 \\
        Staff Review of the Economic and Financial Situation&	  2&	    3286&	 200.82&	3144&	   3215&	  3286&	   3357&	3428 \\
        \bottomrule
    \end{tabular}
    \caption{Summary Statistics for Tokens of Different Sections and Minutes}
    \label{tab:ch2:summary_statis}
\end{table}

% [Codex 修改版] 2026-04-17: 按“保留旧表加注”的策略，明确说明表~\ref{tab:ch2:summary_statis} 不是当前 raw corpus 的主口径。
\noindent\textit{Audit note on Table~\ref{tab:ch2:summary_statis}.}
Table~\ref{tab:ch2:summary_statis} is retained as a \textit{legacy token-statistics snapshot}. Its ``Total Minutes = 123'' row should not be interpreted as the current corpus-wide meeting count. A reproducible raw-data audit of the current repository snapshot identifies 128 unique meeting dates. The discrepancy reflects an older tokenization/statistics export rather than a contradiction in the underlying corpus chronology.


% [CH2-014][Original Start]
% As the maximum length for the major sections are less than 4096 tokens (except the ``Committee Policy Action''), we conduct a cutoff on the text and set the maximum length of the ``Assistant Prompt'' as 4096 to reduce the memory usage.
% [CH2-014][Original End]
Because most major sections are shorter than 4,096 tokens (except ``Committee Policy Action''), we cap the assistant output length at 4,096 tokens to control memory usage.


% Furthermore, the dataset size may not be big enough to conduct fine-tuning. Therefore, we downloaded extra texts as supplementary training materials, including the Transcripts, Green books, Blue books, for re-pre-train the model. The FOMC Transcripts for the corresponding Minutes are the detailed record of the meetings. The FOMC Green books (from 1990 to 2010) contain the detailed analysis of the economic and financial data, and the Blue books (from 1990 to 2010) demonstrate the corresponding economic and financial backgrounds. Table \ref{tab:ch2:summary_supplementary_materials} lists the summary statistics for our supplementary materials.


% \begin{sidewaystable}[htp]
%     \centering
%     \footnotesize
%     \begin{tabular}{c c c c c p{10cm}}
%         \toprule
%         Name&	Start Year&	End Year&	Count&	Average Tokens& \hspace{3.2cm} Description \\
%         \hline
%         Blue Book&	1965&	2010&	424&	15175.19&  The Blue Book, entitled ``Monetary Policy Alternatives'' includes the background and the context on the alternative monetary policy that the committee could consider.\\
%         Green Book&	1964&	2010&	140&	33753.46&  The Green Book, named ``Current Economic and Financial Conditions'', provides in-depth analysis of the U.S. and international economy. \\
%         Beige Book&	2010&	2019&	15&	    21298.2& The Beige Book, entitled ``Summary of Commentary on Current Economic Conditions by Federal Reserve District'', provides information on current economic from different district.\\
%         Teal Book A&	2010&	2024&	69&	71786.20& The Teal Book A, officially subtitled ``Economic and Financial Conditions: Current Situation and Outlook'', provides in-depth analysis and forecasts of the U.S. and international economy (Same as Green Book) and analysis of financial market developments\\
%         Teal Book B&	2010&	2024&	69&	24426.38& The Teal Book B, officially subtitled ``Monetary Policy: Strategies and Alternatives'', provide background and context on monetary policy alternatives that the FOMC could consider at a particular meeting.(Same as Blue Book)\\
%         \makecell[c]{Summary of Economic Projection\\ (SEP)}&	2007&	2019&	20&	7385.30&  The SEP includes a compilation and summary of these projections (which includes individual projections without attribution) is circulated to participants, and a detailed summary of the economic projections\\
%         \bottomrule
%     \end{tabular}
%     \caption{Summary of Supplementary Materials} 
%     \label{tab:ch2:summary_supplementary_materials}
% \end{sidewaystable}

% The Bluebook, officially entitled "Monetary Policy Alternatives," was produced by the staff of the Board of Governors to provide background and context on monetary policy alternatives that the FOMC could consider
% The Greenbook, officially entitled "Current Economic and Financial Conditions," provided in-depth analysis of the U.S. and international economy.
% The Beige Book, known as such originally for its beige cover, is officially entitled "Summary of Commentary on Current Economic Conditions by Federal Reserve District."
% collects anecdotal information on current economic conditions in its District

% Tealbook A, officially subtitled “Economic and Financial Conditions: Current Situation and Outlook,” produced by the staff of the Board of Governors, provides in-depth analysis and forecasts of the U.S. and international economy that were previously covered by Parts 1 and 2 of the Greenbook, and also provides analysis of financial market developments that was previously covered in the Bluebook.
% Tealbook B, officially subtitled “Monetary Policy: Strategies and Alternatives,” is produced by the staff of the Board of Governors to provide background and context on monetary policy alternatives that the FOMC could consider at a particular meeting. Tealbook B includes the same information that had been included in the Bluebook, except that the financial developments content was moved to Tealbook A.
% submit individual economic projections in conjunction with four FOMC meetings a year. A compilation and summary of these projections (which includes individual projections without attribution) is circulated to participants, and a detailed summary of the economic projections




% Bluebook: The Bluebook, officially entitled "Monetary Policy Alternatives," was produced by the staff of the Board of Governors to provide background and context on monetary policy alternatives that the FOMC could consider at an upcoming meeting from 1965 to 2010. It was distributed to FOMC participants during the week before an FOMC meeting--usually one day after the Greenbook.

% Greenbook: The Greenbook, officially entitled "Current Economic and Financial Conditions," provided in-depth analysis of the U.S. and international economy. It was produced by the staff at the Board of Governors and  distributed to FOMC meeting attendees the week before the meeting from 1964 to 2010. For most of its history, the Greenbook was split into two parts plus a supplement:

% Part 1: Current Economic and Financial Conditions: Summary and Outlook
% Part 2: Current Economic and Financial Conditions: Recent Developments
% Supplement: Current Economic and Financial Conditions: Supplemental Information



% Tealbook: Bluebook and Greenbook were merged in June 2010

% Statement on Longer-Run Goals and Monetary Policy Strategy: First released in January 2012, this statement articulates the framework for monetary policy and serves as the foundation for the Committee's policy actions.


% Beige Books: The Beige Book, known as such originally for its beige cover, is officially entitled "Summary of Commentary on Current Economic Conditions by Federal Reserve District." It is produced by Reserve Bank staff and released to the public approximately two weeks prior to each regularly scheduled FOMC meeting.



\subsection*{From FOMC Minutes to Analysis}

% Because the analysis is primarily based on economic and financial market numerical data, it is necessary to extract the relevant data as part of the input prompt. This enables the model to perform logical analysis on the data and generate outputs resembling the style and content of the FOMC Minutes. We employed ``ChatGPT-4o'' as a classifier to label paragraphs from the Minutes using a list of target indicators.
% [Codex 修改版] 2026-02-26: 补充“LLM标注=弱监督”的方法定位与潜在偏差，并提前说明为何剔除 Policy Action（避免决策泄漏）。
Because our prompts are constructed from macroeconomic and financial indicators, we must map unstructured Minutes paragraphs to the corresponding indicator topics. We use an LLM (ChatGPT-4o) as a \textit{weakly supervised classifier} to assign indicator labels to Minutes paragraphs, which enables systematic alignment between numerical inputs and textual targets. Since LLM-based labeling can introduce noise or bias, we mitigate this risk by (i) manually cleaning and standardizing the label ontology and (ii) relabeling with a constrained reference list after cleaning. Importantly, we exclude administrative paragraphs and the ``Committee Policy Action'' content from the \textit{analysis-generation} dataset to avoid leaking realized policy decisions into training, which is crucial for a fair decision-prediction evaluation later in the chapter.



% [CH2-015][Original Start]
% To construct the list of target indicators, we first asked the LLM to label all paragraphs without providing a predefined reference list. The model assigned labels based on its pre-trained knowledge, these are referred to as raw labels. However, several issues were observed with these raw labels, requiring further cleaning and standardization.
% [CH2-015][Original End]
To construct candidate indicator labels, we first asked the LLM to label paragraphs without a predefined reference list. These model-generated tags are treated as \textit{raw labels}. We then identified several quality issues that required cleaning and standardization.


First, some of the raw labels were too general or lacked measurable indicators—for example, labels like ``Housing'', ``International Markets'', ``Swap Arrangements'', and ``Liquidity Facilities''. These categories do not correspond to specific economic indicators.


Second, different labels were often used for the same underlying indicator. For instance, ``PCE Price Index'', ``Consumer Spending'', and ``Personal Consumption Expenditures (PCE)'' all refer to the same economic concept and were consolidated under ``Personal Consumption Expenditures (PCE)''.


% [CH2-016][Original Start]
% Finally, some labels represented subcomponents of broader measures. For example, ``Total Assets of the Federal Reserve'' and ``Total Liabilities of the Federal Reserve'' are both subset of the ``Federal Reserve Balance Sheet''.
% [CH2-016][Original End]
Finally, some labels represented subcomponents of broader constructs. For example, ``Total Assets of the Federal Reserve'' and ``Total Liabilities of the Federal Reserve'' were consolidated under ``Federal Reserve Balance Sheet''.



Therefore, we manually cleaned the \textit{raw labels} and created a list of \textit{reference labels} based on them. We then asked the LLM to relabel the paragraphs using this refined set of reference labels.



% Our raw text dataset consists of 18,957 paragraphs. We removed all paragraphs related to ``Secretary'', ``Attendance'', ``Voting for this action'', and ``Committee Policy Action''. The first three sections contain no analytical content, while the last section only outlines the committee's final decision without providing substantive analysis. As the primary focus of this stage is on analytical content rather than monetary policy decision-making, these exclusions were necessary. After the filtering process, a total of 6,266 paragraphs remained for subsequent labeling.
% [Codex 修改版] 2026-02-26: 强化“过滤规则”的逻辑——既为了分析纯度，也为了避免决策信息泄漏。
Our raw corpus contains 18,957 paragraphs. We removed paragraphs related to ``Secretary'', ``Attendance'', and ``Voting for this action'' because they contain no substantive analytical content. We also removed ``Committee Policy Action'' paragraphs from the analysis-generation dataset, because these paragraphs primarily summarize realized decisions rather than the underlying discussion; excluding them also reduces the risk of decision-information leakage when we later evaluate the model’s ability to infer policy actions from analytical context. After filtering, 6,266 paragraphs remained for indicator labeling.


% [CH2-017][Original Start]
% Next, we instructed the LLM to assign the label ``Non-Core'' to paragraphs that contain no analysis, and ``Other'' when the model was unable to identify a relevant indicator. Since some paragraphs refer to multiple indicators, the model was allowed to extract all applicable indicators and separate them using commas (\textit{,}). This process yielded 5,290 labeled paragraphs with meaningful indicators, which expanded to 5,397 rows across 37 unique labels .
% [CH2-017][Original End]
Next, we introduced two fallback tags: ``Non-Core'' for paragraphs without analytical content, and ``Other'' when no relevant indicator could be mapped. Because some paragraphs refer to multiple indicators, the model was allowed to assign multiple labels separated by commas. This stage produced 5,290 labeled analytical paragraphs, which expanded to 5,397 rows across 37 raw labels.


% [CH2-018][Original Start]
% Subsequently, we manually reviewed the \textit{raw labels}, standardizing and clustering them into 26 new categories to form our final list of \textit{reference labels}. We then asked the model to relabel all 6,266 paragraphs using these \textit{reference labels}, assigning ``Other'' if no relevant indicator could be identified, and ``Non-Core'' if no analysis found in the paragraphs.
% [CH2-018][Original End]
We then manually standardized the raw labels into 26 \textit{reference-label} categories and relabeled all 6,266 filtered paragraphs using this controlled taxonomy. ``Other'' was assigned when no indicator matched, and ``Non-Core'' when the paragraph contained no analytical content.



% [CH2-019][Original Start]
% Finally, we extracted 4,889 analytical paragraphs with a total of 5,781 corresponding standardized labels for our \textbf{main dataset}. We applied the same cleaning process to our \textbf{supplementary dataset}, resulting in 4,464 analytical paragraphs associated with 44,178 standardized labels. Table~\ref{tab:ch2:indicator_count} summarizes the frequency distribution of the reference labels for both datasets. Notably, the supplementary dataset contains significantly more labels than the main dataset. This discrepancy arises because the FOMC Minutes prior to 2009 were less structured, often referencing multiple indicators within a single paragraph. Consequently, for these earlier Minutes, we instructed the model to assign all relevant labels to each paragraph. In contrast, Minutes published after 2009 typically focused on one indicator per paragraph; thus, we asked the model to assign only the single most relevant label in these cases.
% [CH2-019][Original End]
The final \textbf{main dataset} contains 4,887 analytical paragraphs with 5,781 standardized labels. Applying the same procedure to the \textbf{supplementary dataset} yields 4,464 analytical paragraphs with 44,178 labels. Table~\ref{tab:ch2:indicator_count} reports the label distribution. The supplementary set has many more labels because pre-2009 Minutes are less standardized and often mention multiple indicators within a single paragraph; therefore, multi-label assignment is more frequent in that period.




\begin{table}[htbp]
\centering
\begin{tabular}{lcc}
\toprule
\textbf{Indicator} & \textbf{Main Dataset} & \textbf{Supplementary Dataset} \\
\hline
GDP Growth & 623 & 2339 \\
Federal Funds Rate & 614 & 2859 \\
Personal Consumption Expenditures (PCE) & 592 & 3517 \\
Bank Credit to Private Sector & 506 & 1683 \\
Federal Reserve Balance Sheet & 324 & 1381 \\
Consumer Price Index (CPI) & 295 & 2383 \\
Unemployment Rate & 287 & 2276 \\
Business Investment & 232 & 2300 \\
Labour Market & 225 & 1893 \\
Government Purchases & 203 & 1070 \\
Treasury Yields & 200 & 1557 \\
Industrial Production & 164 & 2234 \\
Equity Market Indices & 135 & 1723 \\
Others & 927 & 16901 \\
% Housing Starts & 130 & 2121 \\
% Trade Balance & 115 & 2560 \\
% Exchange Rate & 96 & 2125 \\
% Overnight Rate & 94 & 82 \\
% International Equity Markets & 93 & 349 \\
% Corporate Bond Yields & 82 & 964 \\
% Market Volatility (VIX) & 80 & 514 \\
% Mortgage Rates & 56 & 1586 \\
% Home Prices & 47 & 1067 \\
% Commodity Prices & 45 & 1750 \\
% Money Supply & 39 & 2128 \\
% Bank Capital & 37 & 26 \\
% Consumer Confidence Index & 13 & 1629 \\
\midrule
~\\
Total & 5327 & 44116 \\
\bottomrule
\end{tabular}
\caption{Indicators for Main and Supplementary Datasets}
\label{tab:ch2:indicator_count}
\end{table}




    







\subsection*{Economic and Financial Market Indicators}

% [CH2-020][Original Start]
% After processing the FOMC Minutes, we extracted relevant numerical data to construct the input prompts for our model. Based on our reference labels, we obtained both financial data from WRDS and economic data from the FRED Economic Database.
% [CH2-020][Original End]
After processing the Minutes, we extracted numerical series to build model prompts. Guided by the reference-label taxonomy, we collected macroeconomic series from the FRED database and financial-market series from WRDS and Refinitiv DataScope Select. WRDS provides standardized access to widely used asset-pricing variables and benchmark series, while DataScope Select is used when instrument-level, contract-specific, or higher-frequency market data are required for constructing and validating market variables.


% [CH2-021][Original Start]
% For economic analysis, we collected U.S. macroeconomic data, including GDP growth, Consumer Price Indices (CPIs), the federal funds rate, unemployment rates, export and import figures, government bond yields, and other key indicators. For financial analysis, we included data such as equity market indices, commodity prices, the market volatility index, the Federal Reserve balance sheet, overnight swap rates, and corporate bond yields, among others.
% [CH2-021][Original End]
For macroeconomic coverage, we include GDP growth, Consumer Price Index measures, Federal Funds Rate series, unemployment measures, trade indicators, and Treasury yields. For financial coverage, we include equity indices, commodity prices, volatility indices, Federal Reserve balance-sheet measures, overnight-rate series, and corporate-bond yields.


% [CH2-022][Original Start]
% Our dataset covers a broad range of economic and financial dimensions relevant to monetary policy decision-making. In total, we gathered 102 indicators across 26 categories. Table~\ref{tab:ch2:extracted_data} lists all the indicators along with their corresponding data sources used in the analysis.
% [CH2-022][Original End]
Overall, the dataset spans 102 indicators in 26 categories that are relevant for monetary-policy analysis. Table~\ref{tab:ch2:extracted_data} lists the indicators used in this chapter.

To ensure reproducibility of the DataScope component, we implemented a lightweight Python toolkit (\texttt{datascope\_apis}) that wraps the DataScope Select REST endpoints used in this thesis. The client reads credentials from environment variables (e.g., \texttt{DATASCOPE\_USERNAME}, \texttt{DATASCOPE\_PASSWORD}) with a local \texttt{application.ini} used only as a non-sensitive fallback. Authentication is handled via a cached token with expiry-based refresh and retry/backoff controls. Extractions follow the API's asynchronous job model: requests are submitted to the extraction endpoint, a polling location/job identifier is returned, and the final output is downloaded once the job completes. The toolkit separates request-object construction (template-specific \texttt{*Extractor} classes that serialize to JSON bodies) from HTTP transport (auth/search/extraction/download clients), which supports extension to new templates and field enumerations.

Large-scale queries are I/O-bound from the client perspective and are therefore parallelized by time-window slicing (e.g., monthly/daily partitions) and bounded multithreaded execution (via \texttt{ExtractionImp}/thread executors). Outputs are persisted as compressed \texttt{.csv.gz} files, and a file-integrity validator (\texttt{OutputFileIntegrityValidator}) is used both before and after downloading to identify missing or incomplete files, enabling deterministic reruns and checkpointed recovery. Listing~\ref{code: ch2: datascope intraday example} provides a minimal example of extracting intraday summary data and writing the result to disk.

\begin{lstlisting}[caption = {Minimal DataScope Select intraday extraction (simplified)}, label={code: ch2: datascope intraday example}]
from datetime import date

from src.connection.apis import ExtractionCreator
from src.connection.extraction.enums.extraction_base_enums import IdentifierType
from src.connection.extraction.enums.content_field_names.tick_history.intraday_content_field_names import (
    IntradaySummariesContentFieldNames,
)
from src.connection.utils.condition.condition import TickHistorySummaryInterval
from src.market_data.output_validator import OutputFileIntegrityValidator


identifier = "FFZ4"  # example: 30-day Federal Funds futures contract (RIC)
identifier_type = IdentifierType.Ric

fields = [
    IntradaySummariesContentFieldNames.Close.Bid,
    IntradaySummariesContentFieldNames.Close.Ask,
]

extractor = ExtractionCreator.tick_history(
    contract_id=identifier,
    identifier_type=identifier_type,
    query_start_date=date(2014, 1, 1),
    query_end_date=date(2014, 2, 1),
    content_field_names=fields,
    summary_interval=TickHistorySummaryInterval.FiveMinutes,
)

output_path = "./output/FFZ4_2014-01-01_2014-02-01.csv.gz"
if OutputFileIntegrityValidator.get_missed_files_list(output_path):
    extractor.save_output_file(output_path)
\end{lstlisting}

\begin{center}
\begin{singlespace}    
\begin{scriptsize}
\begin{longtable}{p{0.33\textwidth}|p{0.66\textwidth}}
        \toprule
        \textbf{Category}  & \textbf{Indicator} \\
        \hline
        Labour Market  & All Employees, Total Nonfarm (Thousands of Persons, Seasonally Adjusted) \\
        \hline
        Labour Market  & Job Openings, Total Nonfarm(Level in Thousands, Seasonally Adjusted) \\
        \hline
        Labour Market  & Labor Force Participation Rate(Percent, Seasonally Adjusted) \\
        \hline
        Labour Market  & Total Nonfarm Private Payroll Employment(Persons, Seasonally Adjusted) \\
        \hline
        Money Supply  & Total Monetary Base( currency in circulation plus reserve balances, Billions of Dollars, Not Seasonally Adjusted) \\
        \hline
        Money Supply  & M2SL(Billions of Dollars, Seasonally Adjusted) \\
        \hline
        Unemployment Rate  & Unemployment Rate (Seasonal adjusted) \\
        \hline
        International Equity Markets  & MSCI WORLD \\
        \hline
        International Equity Markets  & MSCI Emerging Markets Index(future) \\
        \hline
        Housing Starts  & Annual Rate for Housing Units Authorized in Permit-Issuing Places in United States (Seasonally Adjusted Total Units [Thousands of Units]) \\
        \hline
        Housing Starts  & New Privately-Owned Housing Units Started, Total Units(Thousands of Units, Seasonally Adjusted Annual Rate) \\
        \hline
        Home Prices  & Rental Vacancy Rate in the United States(Percent, Not Seasonally Adjusted) \\
        \hline
        Home Prices  & All-Transactions House Price Index for the United States( Index 1980Q1=100, Not Seasonally Adjusted) \\
        \hline
        Home Prices  & S\&P CoreLogic Case-Shiller U.S. National Home Price Index (Index Jan 2000=100, Seasonally Adjusted) \\
        \hline
        Federal Reserve Balance Sheet  & Balance Sheet,Total Assets(Millions of U.S. Dollars, Not Seasonally Adjusted) \\
        \hline
        Federal Reserve Balance Sheet  & Balance Sheet, Mortgage-Backed Securities(Millions of U.S. Dollars, Not Seasonally Adjusted) \\
        \hline
        Federal Reserve Balance Sheet  & Reserve Balances with Federal Reserve Banks(Billions of U.S. Dollars, Not Seasonally Adjusted) \\
        \hline
        Federal Reserve Balance Sheet  & Summary Of SOMA holdings \\
        \hline
        Exchange Rate  & U.S. Dollars to Euro Spot Exchange Rate(U.S. Dollars to One Euro, Not Seasonally Adjusted) \\
        \hline
        Exchange Rate  & U.S. Dollars to U.K. Pound Sterling Spot Exchange Rate(U.S. Dollars to One U.K. Pound Sterling, Not Seasonally Adjusted) \\
        \hline
        Exchange Rate  & Japanese Yen to U.S. Dollar Spot Exchange Rate(Japanese Yen to One U.S. Dollar, Not Seasonally Adjusted) \\
        \hline
        Exchange Rate  & Nominal Broad U.S. Dollar Index(Index Jan 2006=100, Not Seasonally Adjusted) \\
        \hline
        Exchange Rate  & Chinese Yuan Renminbi to U.S. Dollar Spot Exchange Rate(Chinese Yuan Renminbi to One U.S. Dollar, Not Seasonally Adjusted) \\
        \hline
        Consumer Confidence Index  & Advance Monthly Sales for Retail and Food Services(Seasonally Adjusted Sales - Monthly [Millions of Dollars]) \\
        \hline
        Consumer Confidence Index  & Consumer Sentiment(Index 1966Q1=100, Not Seasonally Adjusted) \\
        \hline
        Corporate Bond Yields  & ICE BofA 1-3 Year US Corporate Index Effective Yield(Percent, Not Seasonally Adjusted) \\
        \hline
        Corporate Bond Yields  & ICE BofA CCC \& Lower US High Yield Index Effective Yield(Percent, Not Seasonally Adjusted) \\
        \hline
        Corporate Bond Yields  & ICE BofA AAA US Corporate Index Effective Yield(Percent, Not Seasonally Adjusted) \\
        \hline
        Corporate Bond Yields  & ICE BofA 7-10 Year US Corporate Index Effective Yield(Percent, Not Seasonally Adjusted) \\
        \hline
        Corporate Bond Yields  & ICE BofA BBB US Corporate Index Effective Yield (Percent, Not Seasonally Adjusted) \\
        \hline
        Corporate Bond Yields  & ICE BofA 5-7 Year US Corporate Index Effective Yield(Percent, Not Seasonally Adjusted) \\
        \hline
        Corporate Bond Yields  & ICE BofA 3-5 Year US Corporate Index Effective Yield(Percent, Not Seasonally Adjusted) \\
        \hline
        Corporate Bond Yields  & ICE BofA B \& Lower US Emerging Markets Liquid Corporate Plus Index Effective Yield(Percent, Not Seasonally Adjusted) \\
        \hline
        Corporate Bond Yields  & ICE BofA 15+ Year US Corporate Index Effective Yield(Percent, Not Seasonally Adjusted) \\
        \hline
        Corporate Bond Yields  & ICE BofA BB US High Yield Index Effective Yield (Percent, Not Seasonally Adjusted) \\
        \hline
        Corporate Bond Yields  & ICE BofA AA US Corporate Index Effective Yield(Percent, Not Seasonally Adjusted) \\
        \hline
        Corporate Bond Yields  & ICE BofA 10-15 Year US Corporate Index Effective Yield(Percent, Not Seasonally Adjusted) \\
        \hline
        Corporate Bond Yields  & Moody's Seasoned Baa Corporate Bond Yield(Percent, Not Seasonally Adjusted) \\
        \hline
        Corporate Bond Yields  & Moody's Seasoned Aaa Corporate Bond Yield(Percent, Not Seasonally Adjusted) \\
        \hline
        Commodity Prices  & Producer Price Index by Commodity, Fuels and Related Products and Power (Index 1982=100, Not Seasonally Adjusted) \\
        \hline
        Commodity Prices  & Producer Price Index by Commodity, Farm Products(Index 1982=100, Not Seasonally Adjusted) \\
        \hline
        Commodity Prices  & Producer Price Index by Commodity, All Commodities(Index 1982=100, Not Seasonally Adjusted) \\
        \hline
        Commodity Prices  & Crude Oil Prices(West Texas Intermediate, Dollars per Barrel, Not Seasonally Adjusted) \\
        \hline
        Commodity Prices  & Producer Price Index by Commodity, Metals and Metal Products(Index 1982=100, Not Seasonally Adjusted) \\
        \hline
        Commodity Prices  & Crude Oil Prices(Brent,Dollars per Barrel, Not Seasonally Adjusted) \\
        \hline
        GDP Growth  & Gross domestic product ([Millions of dollars] Seasonally adjusted at annual rates) \\
        \hline
        Mortgage Rates  & 30-Year Fixed Rate Mortgage Average in the United States(Percent, Not Seasonally Adjusted, Weekly, Ending Thursday) \\
        \hline
        Mortgage Rates  & 15-Year Fixed Rate Mortgage Average in the United States(Percent, Not Seasonally Adjusted) \\
        \hline
        Equity Market Indices  & Dow Jones Industrial Average \\
        \hline
        Equity Market Indices  & NASDAQ Composite Index \\
        \hline
        Equity Market Indices  & S\&P 500 Index \\
        \hline
        Industrial Production  & industrial\_production(total manufacturing, Index 2017=100, Seasonally Adjusted) \\
        \hline
        Industrial Production  & Industrial Production( Manufacturing, Durable Goods, Semiconductor and Other Electronic Component ) \\
        \hline
        Industrial Production  & Industrial Production( Mining) \\
        \hline
        Industrial Production  & Industrial Production( Manufacturing, Nondurable Goods, Food, Beverage, and Tobacco) \\
        \hline
        Industrial Production  & Industrial Production( Utilities, Electric and Gas Utilities) \\
        \hline
        Industrial Production  & Industrial Production(Manufacturing, Durable Goods,  Motor Vehicles and Parts) \\
        \hline
        Industrial Production  & industrial\_production(total\_index, Index 2017=100, Seasonally Adjusted) \\
        \hline
        Business Investment  & Real Private Nonresidential Fixed Investment(Billions of Chained 2017 Dollars, Seasonally Adjusted Annual Rate) \\
        \hline
        Business Investment  & Real Gross Private Domestic Investment(Billions of Chained 2017 Dollars, Seasonally Adjusted Annual Rate) \\
        \hline
        Business Investment  & Commercial Paper Outstanding(Billions of Dollars, Seasonally Adjusted) \\
        \hline
        Personal Consumption Expenditures (PCE)  & Personal Consumption Expenditures Index(Index 2017=100, Seasonally Adjusted) \\
        \hline
        Personal Consumption Expenditures (PCE)  & Personal consumption expenditures by major function (Index numbers, 2017=100; quarters and months are seasonally adjusted) \\
        \hline
        Personal Consumption Expenditures (PCE)  & Personal consumption expenditures by major function(Millions of dollars; quarters and months are seasonally adjusted at annual rates) \\
        \hline
        Market Volatility (VIX)  & CBOE S\&P 500 3-Month Volatility Index(Index, Daily, Not Seasonally Adjusted) \\
        \hline
        Market Volatility (VIX)  & CBOE Gold ETF Volatility Index(Index, Not Seasonally Adjusted) \\
        \hline
        Market Volatility (VIX)  & CBOE Crude Oil ETF Volatility Index(Index, Not Seasonally Adjusted) \\
        \hline
        Market Volatility (VIX)  & CBOE DJIA Volatility Index (Index, Daily, Not Seasonally Adjusted) \\
        \hline
        Market Volatility (VIX)  & CBOE NASDAQ 100 Volatility Index(Index, Daily, Not Seasonally Adjusted) \\
        \hline
        Market Volatility (VIX)  & CBOE Volatility Index(Index, Not Seasonally Adjusted) \\
        \hline
        Market Volatility (VIX)  & CBOE Russell 2000 Volatility Index(Index, Daily, Not Seasonally Adjusted) \\
        \hline
        Trade Balance  & Foreign Transactions in the National Income and Product Accounts (Millions of dollars Seasonally adjusted at annual rates) \\
        \hline
        Federal Funds Rate  & Federal Funds Effective Rate(Percent, Not Seasonally Adjusted) \\
        \hline
        Federal Funds Rate  & Federal Funds Target Range - Upper Limit (Percent, Not Seasonally Adjusted) \\
        \hline
        Bank Capital  & Bank Total Assets, All Commercial Banks(Billions of U.S. Dollars, Seasonally Adjusted) \\
        \hline
        Bank Capital  & Bank Capital to Total Assets for United States(Percent, Not Seasonally Adjusted) \\
        \hline
        Overnight Rate  & Secured Overnight Financing Rate (Percent, Not Seasonally Adjusted) \\
        \hline
        Government Purchases  & Government current expenditures, National defense(Billions of Dollars, Not Seasonally Adjusted) \\
        \hline
        Government Purchases  & Real Government Consumption Expenditures and Gross Investment (Billions of Chained 2017 Dollars, Seasonally Adjusted Annual Rate) \\
        \hline
        Government Purchases  & Federal Government,Current Expenditures(Billions of Dollars, Seasonally Adjusted Annual Rate) \\
        \hline
        Consumer Price Index (CPI)  & Consumer Price Index for All Urban Consumers (CPI-U, seasonal adjusted) \\
        \hline
        Bank Credit to Private Sector  & SLOOS,Net Percentage of Domestic Banks Tightening Standards for Credit Card Loans(Percent, Not Seasonally Adjusted) \\
        \hline
        Bank Credit to Private Sector  & Total Consumer Credit Owned and Securitized(Millions of Dollars, Seasonally Adjusted) \\
        \hline
        Bank Credit to Private Sector  & SLOOS,Net Percentage of Domestic Banks Tightening Standards for Commercial and Industrial Loans to Large and Middle-Market Firms(Percent, Not Seasonally Adjusted) \\
        \hline
        Bank Credit to Private Sector  & SLOOS, Net Percentage of Domestic Banks Tightening Standards for Auto Loans(Percent, Not Seasonally Adjusted) \\
        \hline
        Bank Credit to Private Sector  & Personal Saving Rate(Percent, Seasonally Adjusted Annual Rate) \\
        \hline
        Bank Credit to Private Sector  & SLOOS, Net Percentage of Domestic Banks Reporting Increased Willingness to Make Consumer Installment Loans(Percent, Not Seasonally Adjusted) \\
        \hline
        Bank Credit to Private Sector  & SLOOS,Net Percentage of Domestic Banks Tightening Standards for Commercial and Industrial Loans to Small Firms(Percent, Not Seasonally Adjusted) \\
        \hline
        Bank Credit to Private Sector  & Percent Change of Total Consumer Credit(Percent Change at Annual Rate, Seasonally Adjusted Annual Rate) \\
        \hline
        Bank Credit to Private Sector  & KBW Nasdaq Bank Index \\
        \hline
        Bank Credit to Private Sector  & SLOOS, Net Percentage of Domestic Banks Tightening Standards for Commercial Real Estate Loans with Construction and Land Development Purposes(Percent, Not Seasonally Adjusted) \\
        \hline
        Treasury Yields  & Market Yield on U.S. Treasury Securities at 3-Month Constant Maturity(Percent, Not Seasonally Adjusted) \\
        \hline
        Treasury Yields  & Market Yield on U.S. Treasury Securities at 1-Month Constant Maturity(Percent, Not Seasonally Adjusted) \\
        \hline
        Treasury Yields  & Market Yield on U.S. Treasury Securities at 7-Year Constant Maturity(Percent, Not Seasonally Adjusted) \\
        \hline
        Treasury Yields  & Market Yield on U.S. Treasury Securities at 10-Year Constant Maturity(Percent, Not Seasonally Adjusted) \\
        \hline
        Treasury Yields  & Market Yield on U.S. Treasury Securities at 3-Year Constant Maturity(Percent, Not Seasonally Adjusted) \\
        \hline
        Treasury Yields  & Market Yield on U.S. Treasury Securities at 30-Year Constant Maturity(Percent, Not Seasonally Adjusted) \\
        \hline
        Treasury Yields  & Market Yield on U.S. Treasury Securities at 5-Year Constant Maturity(Percent, Not Seasonally Adjusted) \\
        \hline
        Treasury Yields  & Market Yield on U.S. Treasury Securities at 2-Year Constant Maturity(Percent, Not Seasonally Adjusted) \\
        \hline
        Treasury Yields  & Market Yield on U.S. Treasury Securities at 1-Year Constant Maturity(Percent, Not Seasonally Adjusted) \\
        \hline
        Treasury Yields  & Market Yield on U.S. Treasury Securities at 6-Month Constant Maturity(Percent, Not Seasonally Adjusted) \\
        \hline
        Treasury Yields  & Market Yield on U.S. Treasury Securities at 20-Year Constant Maturity(Percent, Not Seasonally Adjusted) \\
        \bottomrule   
        \caption{Economic and Financial Market Indicators}
        \label{tab:ch2:extracted_data}   
\end{longtable}
\end{scriptsize}
\end{singlespace}
\end{center}



\subsection*{Constructing Reasoning Process}

% [CH2-023][Original Start]
% Our next step is to establish a logical link between the analysis found in the FOMC Minutes and the corresponding numerical data. The FOMC Minutes provide high-level summaries intended to support policy decisions, often lacking detailed and explicit reasoning processes. As a result, large language models (LLMs) cannot infer these summaries directly from the raw data. To train the model to generate logical and reliable conclusions, we simulate the underlying reasoning by applying model distillation.
% [CH2-023][Original End]
The next step is to link numerical indicators to the analytical narratives in FOMC Minutes. Minutes typically provide high-level conclusions rather than explicit reasoning chains; therefore, direct mapping from raw data to final narrative is under-specified. We address this by distilling intermediate reasoning traces.

% [CH2-024][Original Start]
% Model distillation (\cite{hinton2015distilling}) involves training a smaller, less powerful model using data generated by a larger, more capable model. Rather than exposing the smaller model to a wide range of general data, distillation provides focused, high-quality data that encapsulates expert-level reasoning. For example, the 8-billion-parameter LLaMA 3 model achieves comparable accuracy on Math benchmarks to OpenAI's ChatGPT-o1-mini, which has around 100 billion parameters, after fine-tuned with distilled, high-quality datasets.
% [CH2-024][Original End]
Model distillation (\cite{hinton2015distilling}) trains a smaller model using outputs from a stronger teacher model. Instead of broad general-domain supervision, distillation supplies focused, task-relevant traces that encode higher-quality reasoning patterns.


% [CH2-025][Original Start]
% More specifically, we applied the ``DeepSeek-R1'' as our teacher model to generate the reasoning process. Table~\ref{tab:ch2:prompt_for_data_distillation} demonstrates the prompt we used for data distillation. We used the \texttt{meeting date}, \texttt{section name}, \texttt{topic} of the indicators, \texttt{table} contained the indicators' data, and the \texttt{reference excerpt} from the FOMC Minutes to construct the prompt. Noticeably, the \texttt{reference excerpt} provided a reference to the model to ensure the relevance of the output to the FOMC Minutes.
% [CH2-025][Original End]
In implementation, we use DeepSeek-R1 as the teacher model for reasoning-trace generation. As shown in Table~\ref{tab:ch2:prompt_for_data_distillation}, each prompt includes the meeting date, target section style, indicator topic, indicator tables, and a reference excerpt from FOMC Minutes. The excerpt anchors style and topical relevance during \textit{training-data construction}. Because that excerpt is closely tied to the target text, it should not be interpreted as evidence that the downstream train/eval/test splits are leakage-free.


\begin{table}[htbp]
    \centering
    \footnotesize
    \begin{tabular}{p{0.98\textwidth}}
    \toprule
    \textbf{Prompt for Data Distillation} \\
    \midrule
% [CH2-026][Original Start]
% You are an economist preparing briefing notes for the Federal Open Market Committee (FOMC) meeting on \texttt{**\{meeting\_date\}**}.  
% Your task is to analyze recent trends in \texttt{**\{topic\}**} and related indicators, writing in the tone and style of the \texttt{**\{section\_style\}**} section of the FOMC minutes. \\
% [CH2-026][Original End]
You are an economist preparing briefing notes for the Federal Open Market Committee (FOMC) meeting on \texttt{**\{meeting\_date\}**}.  
Your task is to analyze recent developments in \texttt{**\{topic\}**} and related indicators, using the tone and structure of the \texttt{**\{section\_style\}**} section of the FOMC Minutes. \\

Use neutral, data-driven, and measured language throughout your response. \\

Your analysis must be based on the following indicators: \texttt{**\{emphasized\_label\}**}, using the accompanying data tables from the previous two years: \\

~\\
\texttt{\{table\_str\}} \\
~\\


Model your tone, structure, and analytical framing on the following excerpt from the FOMC Minutes:\\

~\\
\texttt{\{reference\_excerpt\}} \\
~\\

Keep your response focused, analytically rigorous, and under 4,096 tokens. \\
    \bottomrule
    \end{tabular}
    \caption{Prompt for Data Distillation}
    \label{tab:ch2:prompt_for_data_distillation}
\end{table}

% [CH2-027][Original Start]
% Finally, we collected 4,889 messages of ``Question'', ``Reasoning'', and ``Answer'' pairs as our dataset for fine-tuning.
% [CH2-027][Original End]
Finally, this procedure produced 4,887 ``Question--Reasoning--Answer'' training samples for fine-tuning.




\subsection*{Synthetic Text Generation}

% [CH2-028][Original Start]
% To assess the simulation capabilities of large language models (LLMs) in structured text, we conducted additional fine-tuning, explicitly instructing the model to produce analyses aligned with the style and structure of actual FOMC Minutes.
% [CH2-028][Original End]
To evaluate structured-text simulation, we add a downstream alignment stage that trains the model to rewrite analysis into authentic FOMC-Minutes style.


% %%%%%%%%%%
% Plan: 
% Data -> analysis -> summary -> FOMC sections

% individual indicator data -> individual analysis -> combined analysis -> summary -> FOMC section



% [CH2-029][Original Start]
% To achieve this, we designed a pipeline that accepts individual indicator data as input and produces structured text aligned with sections from actual FOMC Minutes as output. The process comprises three main steps: analysis generation, summary creation, and section rewriting. First, the model analyzes individual economic or financial indicators and generates analytical paragraphs for each. Next, analyses for several related indicators are combined, and a concise summary is created to serve as input for the subsequent step. Finally, the model rewrites this summary into text that closely follows the format and stylistic conventions of the official FOMC Minutes.
% [CH2-029][Original End]
The pipeline accepts indicator data and outputs section-aligned Minutes text in three steps: (i) indicator-level analysis generation, (ii) consolidation into a concise analytical summary, and (iii) section-specific rewriting into institutional FOMC style.


% [CH2-030][Original Start]
% Since the model had already demonstrated robust proficiency in economic and financial analysis from last stage, the primary objective at this stage was to teach it to generate structured text precisely aligned with specific formatting guidelines.
% [CH2-030][Original End]
Because the model already demonstrates baseline analytical capability from the prior stage, the objective here is style and structure alignment rather than new domain learning.



% [CH2-031][Original Start]
% To facilitate this, we reorganized the FOMC Minutes data according to their respective section titles, using the model-generated summaries as inputs and the actual text from corresponding sections of the Minutes as target outputs. The model was then instructed to rewrite the provided analytical summaries to match the style and format of the authentic FOMC sections. Table~\ref{tab:ch2:fomc_summary_prompt} illustrates the prompt structure and corresponding expected outputs for this synthetic FOMC-style text generation task.
% [CH2-031][Original End]
Accordingly, we reorganize Minutes data by section title, use model-generated summaries as inputs, and use the corresponding official section text as targets. The model is trained to rewrite summaries into section-consistent FOMC language. Table~\ref{tab:ch2:fomc_summary_prompt} shows the prompt template.


\begin{table}[htbp]
\centering
\footnotesize
\begin{tabular}{p{0.96\textwidth}}
\toprule
\textbf{User Prompt: FOMC Minutes-Style Economic Summary Generation} \\
\midrule
You are a senior economist at the Federal Reserve responsible for drafting internal briefing materials that may be incorporated into the official \textbf{FOMC Minutes}. \\
Your task is to transform the following economic or financial analysis into a formal, institutionally appropriate paragraph aligned with the tone and structure of the \textbf{\{section\_name\}} section in the FOMC Minutes for the \textbf{\{meeting\_date\}} meeting. \\[1ex]
~\\
\textbf{\#\#\# You will be provided with:}\\
    - \textbf{\textless Raw Analysis\textgreater}: A semi-structured or informal analytical paragraph that may include market commentary, data points, or economic observations. \\
    - \textbf{\textless Target Section\textgreater}: The FOMC Minutes section in which the rewritten paragraph would appear (e.g., \textit{“Economic Outlook,” “Financial Markets,” “Monetary Policy Considerations”}). \\
~\\
\textbf{\#\#\# Your Objective:} \\
   - Rewrite the content in the formal, precise, and policy-grounded style used in FOMC Minutes. \\
   -  Use an \textbf{objective, technical, and neutral tone} — avoid informal language, first-person voice, or speculative remarks. \\
   - Where appropriate, incorporate phrases typical of FOMC drafting style, such as: \\
   \textbullet\ ``Participants noted that...'' \\
    \textbullet\ ``The Committee observed that...'' \\
    \textbullet\ ``Recent indicators suggested...'' \\
~\\
\textbf{\#\#\# Output Requirements:} \\
- Generate a \textbf{single, coherent paragraph} between \textbf{100–200 words} in FOMC Minutes style. \\
- \textbf{Do not merely summarize} — restructure and elevate the input to meet institutional standards. \\
- Avoid lists or bullet points; use complete sentences and formal syntax. \\
~\\
\textbf{\#\#\# FOMC Minutes Section Name:} \\[1ex]
\inputtag{Target Section} \\
\texttt{\{section\_name\}} \\
\inputtag{/Target Section} \\[1ex]
~\\
\textbf{\#\#\# Input Content:} \\
\inputtag{Raw Analysis} \\
\texttt{\{raw\_analysis\}} \\
\inputtag{/Raw Analysis} \\
\bottomrule
\toprule
\textbf{Target Output:} \\
\midrule
\texttt{\{Corresponding FOMC Section\}} \\
\bottomrule
\end{tabular}
\caption{Prompt for FOMC-Style Text Generation}
\label{tab:ch2:fomc_summary_prompt}
\end{table}


% [CH2-032][Original Start]
% After that, we need to construct the reasoning process ($c_i$) in the model's target response. As we have the target answer ($a_i$), here we applied a standard model distillation pipeline. Firstly, we use the prompt in Table~\ref{tab:ch2:fomc_summary_prompt} in prompting a powerful reasoning model (we used DeepSeek-R1 in our experiment). For each prompt, we generated 5 candidate responses with a high hyperparameters (temperature = 0.6 and top\_p = 0.9). Then, we calculate the Cosine Similarity and the BertScore between each answer part from the responses and the FOMC Minutes' section. For each section, we choose two responses with the highest total score, and use the corresponding reasoning part as our reasoning process ($c_i$).
%
% Next, we constructed the reasoning process ($c_i$) in the model's target response. Given the target answer ($a_i$), we adopted a standard model distillation pipeline. Specifically, we used the prompt shown in Table~\ref{tab:ch2:fomc_summary_prompt} to query a powerful reasoning model, (DeepSeek-R1 in our experiment). For each prompt, we generated five candidate responses using high sampling parameters (temperature = 0.6; top\_p = 0.9). We then evaluated the answer portion of each response by computing both the sentence-level cosine similarity and BERTScore (\cite{zhang2019bertscore}) against the corresponding section of the FOMC Minutes. For each section, we selected the two responses with the highest total scores, and extracted their reasoning components as the final reasoning processes ($c_i$).
% [CH2-032][Original End]
We then construct the reasoning component ($c_i$) for each target response. Given target answer text ($a_i$), we query a stronger reasoning model (DeepSeek-R1) with the template in Table~\ref{tab:ch2:fomc_summary_prompt}. For each prompt, we generate five candidates with high-sampling settings (temperature = 0.6, top\_p = 0.9), score candidate answers against the reference section using cosine similarity and BERTScore (\cite{zhang2019bertscore}), and keep the two highest-scoring traces. Their reasoning segments are used as distilled $c_i$ targets.



% Finally, we collected 718 sections from 128 FOMC Minutes of the FOMC meeting from Jan 2019 to Jan 2025. Table~\ref{tab:ch2:section_name_distribution} listed the most commonly used section names. After applying the prompt template in Table~\ref{tab:ch2:fomc_summary_prompt}, we constructed 1,392 corresponding Q\&A pair questions with two different reasoning process for each section. 
%
% Finally, we collected 718 sections from 128 FOMC Minutes spanning meetings held between January 2019 and January 2025. Table~\ref{tab:ch2:section_name_distribution} presents the most commonly used section titles. By applying the prompt template shown in Table~\ref{tab:ch2:fomc_summary_prompt}, we constructed 1,392 corresponding question-answer (Q\&A) pairs, each accompanied by two distinct reasoning processes for every section.
% [Codex 修改版] 2026-02-26: 修正明显不可能的时间范围（2019--2025 不可能有 128 次 Minutes），并合并重复段落。
Finally, we collected 718 section-level samples from the post-2009 Minutes corpus used in this study (spanning meetings between January 2009 and January 2025). Table~\ref{tab:ch2:section_name_distribution} presents the most common section titles. By applying the prompt template shown in Table~\ref{tab:ch2:fomc_summary_prompt}, we constructed 1,392 corresponding Q\&A pairs, each accompanied by two distinct distilled reasoning processes for the same target section.



\begin{table}[ht]
    \centering
    \footnotesize
    \begin{tabular}{p{0.65\textwidth} r}
        \toprule
        \textbf{Section Name} & \textbf{Count} \\
        \midrule
        Staff Review of the Financial Situation & 125 \\
        Staff Economic Outlook & 124 \\
        Staff Review of the Economic Situation & 124 \\
        Participants' Views on Current Conditions and the Economic Outlook & 116 \\
        Developments in Financial Markets and Open Market Operations & 70 \\
        Developments in Financial Markets and the Federal Reserve's Balance Sheet & 45 \\
        \bottomrule
    \end{tabular}
    \caption{Distribution of Major Section Names in FOMC Minutes}
    \label{tab:ch2:section_name_distribution}
\end{table}

% [CH2-033][Original Start]
% The dataset is relative small for model fine-tuning. Therefore, we used the Minutes before 2009 to boost our data. However, these Minutes have no clear section titles. We need create the section name for each paragraph and then construct the section. We used the section names from Table~\ref{tab:ch2:section_name_distribution} in the construction and designed a prompt (Table~\ref{tab:ch2:section_name_prompt}) asking the LLM to assign the section name to each paragraph.
% [CH2-033][Original End]
The post-2009 section-aligned sample is relatively small for fine-tuning. To expand coverage, we incorporate pre-2009 Minutes, which do not provide stable section headers. We therefore infer section labels paragraph by paragraph using the section taxonomy in Table~\ref{tab:ch2:section_name_distribution} and the prompt in Table~\ref{tab:ch2:section_name_prompt}.


\begin{table}[htbp]
\centering
\footnotesize
\begin{tabular}{p{0.96\textwidth}}
\toprule
\textbf{User Prompt: Create Section Name} \\
\hline
% [CH2-034][Original Start]
% You are a policy editor at the Federal Reserve tasked with organizing the content of the **FOMC Minutes**. \\
% [CH2-034][Original End]
You are a policy editor at the Federal Reserve responsible for organizing **FOMC Minutes** content. \\

Given a paragraph from the FOMC Minutes, assign it to the **single most appropriate section title** from the official list below: \\

**Available Section Titles**: \\

\texttt{**\{Section Names\}**}

--- \\
~\\
\#\#\# Instructions: \\

- Read the paragraph carefully and select **one** section title that best reflects its content and purpose. \\
- Choose the **closest thematic match** based on subject matter, phrasing, and policy focus. \\
- If multiple titles appear plausible, select the one that best reflects **official FOMC categorization**.\\
- Your output must follow this format:   \\
  `\textbackslash boxed\{\{section\_title\}\}` \\

--- \\
~\\
\#\#\# FOMC Paragraph: \\
\inputtag{Paragraph}\\
\texttt{\{fomc\_paragraph\}}  \\
\inputtag{/Paragraph}\\

--- \\
~\\
\#\#\# Output: \\
\textbackslash \\boxed\{\{Your selected section title from the list above\}\} \\

Also provide a brief rationale (1--2 sentences) below the boxed output. \\

\bottomrule
\end{tabular}
\caption{Prompt for Section Name Creation}
\label{tab:ch2:section_name_prompt}
\end{table}










% \subsection*{Monetary Policy Decision Making}
% The final section in the FOMC Minutes is the committee's action on the monetary policy. Therefore, to explore the capability of LLMs in predicting government monetary policy—specifically in influencing the interest rate market, we ask the model to specialize in making decisions on the Effective Federal Funds Rate (EFFR). As the model was trained to conduct logical analyses based on provided data, the objective at this stage was to enable it to make informed decisions grounded in that analysis.

% To this end, we first extracted historical EFFR changes from the Federal Reserve Bank of New York. We then designed a tailored prompt to guide the model in predicting government interest rate decisions. Table~\ref{tab:ch2:prompt_for_decision_prediction} presents the template used for decision prediction prompts.


% \begin{table}[htp]
%     \centering
%     \footnotesize
%     \begin{tabular}{p{0.98\textwidth}}
%     \toprule
%     \textbf{Prompt for Government Decision Prediction} \\
%     \midrule
% You are a voting member of the Federal Open Market Committee (FOMC), responsible for setting U.S. monetary policy. 
% Below is an analysis of U.S. economic and financial conditions ahead of the FOMC meeting on \texttt{**\{meeting\_date\}**}.  \\
% The current Effective Federal Funds Rate is \texttt{**\{current\_rate\}\%** }. \\
% ~\\
% Your assignment:  
% 1. Review the analysis and summarize the key economic and financial factors influencing the decision.  
% 2. Provide a step-by-step explanation of your policy recommendation for this meeting.  
% 3. Conclude with your \textbf{vote}. Select only one of the following:  \\
% \hspace*{0.5cm} - Raise by 100 basis points  \\
% \hspace*{0.5cm} - Raise by 75 basis points  \\
% \hspace*{0.5cm} - Raise by 50 basis points  \\
% \hspace*{0.5cm} - Raise by 25 basis points  \\
% \hspace*{0.5cm} - No change  \\
% \hspace*{0.5cm} - Cut by 25 basis points  \\
% \hspace*{0.5cm} - Cut by 50 basis points  \\
% \hspace*{0.5cm} - Cut by 75 basis points  \\
% \hspace*{0.5cm} - Cut by 100 basis points  \\
% ~\\
% Guidelines:  \\
% - Stay professional, structured, and objective.  \\
% - Do not speculate beyond the information provided. \\  
% ~\\
% \textbf{---} \\
% \textbf{Current Economic and Financial Market Analysis} \\
% \texttt{\{current\_analysis\}} \\
% \textbf{---} \\
% Please structure your response as follows:  \\
% \textbf{1. Summary of Key Economic Factors} (\(\leq 150\) words)  \\
% \textbf{2. Step-by-Step Policy Reasoning}  \\
% \hspace*{0.5cm} - Analyze the data  \\
% \hspace*{0.5cm} - Consider all policy alternatives  \\
% \hspace*{0.5cm} - Justify the your final choice \\
% ~\\
% \textbf{3. Final Vote}  \\
% State your choice as: \texttt{**\textbackslash boxed\{\{your choice\}\}**} (e.g., \texttt{\textbackslash boxed\{\{No change\}\}}). \\
%     \bottomrule
%     \toprule
% \textbf{Target Output:} \\
% \midrule
% \texttt{\{Federal Fund Rate Change\}} \\
% \bottomrule
%     \end{tabular}
%     \caption{Prompt for Government Decision Prediction}
%     \label{tab:ch2:prompt_for_decision_prediction}
% \end{table}


% % Generate many reasoning traces for the prompt.
% % Use a verifier to remove reasoning traces with the wrong final answer.
% % For each remaining trace, extract the set of equations appearing in it. Deduplicate the traces so that each one has a different set of equations. Add those to the dataset.











% \subsection*{Reasoning and Analysis Quality Cleaning}



% Before we move the model fine-tuning, we conduct a filtering to the dataset. As the ``teacher model'' may generate low quality reasoning and answer, we further evaluated both the logical between the data and the reasoning process and the relationship between the data and the answer. Here, we applied ``DeepSeek-R1'' as our evaluator and designed prompts to ask the model scoring the logical processes. A higher score indicate a higher quality output.


% First of all, to ensure that the model generate a reasoning process based on the given data, we evaluate the logical between the data and the reasoning process. We listed 7 aspects as our criteria in the evaluation. 


\subsection*{Dataset Allocation}


% [CH2-035][Original Start]
% Finally, we divided the entire dataset into three parts: Training data, Evaluation data, and Test data. The full dataset comprises 4,889 question-answer (Q\&A) pairs, covering the FOMC Minutes from Jan 2009 to Jan 2025. This dataset was initially divided into 80\% for training, 10\% for evaluation, and 10\% for testing. Within the training subset, 80\% was allocated to the Supervised Fine-Tuning (SFT) step, while the remaining 20\% was reserved for Group Relative Policy Optimization (GRPO). Compared with SFT, GRPO is significantly more computationally intensive; therefore, only a smaller sample of the training data was used for this step. Additionally, to mitigate the risk of GRPO training overwriting knowledge or stylistic features learned during SFT, 10\% of the SFT training data was also included in the GRPO training set. Table~\ref{tab:ch2:dataset_partition_stage1} provides a summary of the dataset partitioning for the training.
% [CH2-035][Original End]
Finally, the canonical Chapter 2 pipeline assigns meetings to training, evaluation, and test subsets at the \textit{meeting} level rather than at the sample level. For the analysis and Minutes-style alignment tasks, we first restrict attention to post-2009 meetings, obtaining 128 unique meetings shared across the two datasets. These meetings are sorted, shuffled with a fixed seed of 42, and split into 102 training meetings, 13 evaluation meetings, and 13 test meetings. All samples derived from a given meeting inherit the same split assignment.

Within the analysis training partition, we use all training rows for SFT. For GRPO, we deterministically sample 20\% of the analysis-training rows as the primary GRPO set and mix in an additional 10\% SFT replay subset to reduce forgetting of format and domain alignment learned during SFT. By contrast, the decision branch starts from the deduplicated union of the current decision GRPO prompts. After collapsing repeated prompts to one canonical record per meeting, the decision pool contains 237 unique meetings, which are then split with the same fixed-random meeting-level procedure into 189 training meetings, 24 evaluation meetings, and 24 test meetings.

\begin{table}[htbp]
\centering
\footnotesize
\begin{tabular}{l l r r r}
\toprule
\textbf{Dataset Family} & \textbf{Population Definition} & \textbf{Train} & \textbf{Eval} & \textbf{Test} \\
\midrule
Analysis SFT / Minutes alignment & Post-2009 meetings shared across the analysis and Minutes datasets & 102 & 13 & 13 \\
Decision GRPO & Deduplicated union of decision prompts, one canonical prompt per meeting & 189 & 24 & 24 \\
\bottomrule
\end{tabular}
\caption{Canonical Meeting-Level Split Counts for the Chapter 2 v2 Pipeline}
\label{tab:ch2:dataset_partition_stage1}
\end{table}

Because the split is random rather than chronological, the resulting experiments should be interpreted as held-out meeting validation under a reproducible meeting-level partition. They do \textit{not} constitute a time-series holdout or a forward-only forecasting design. The exact meeting memberships and row-level manifests are saved by the dataset builder under \texttt{dataset/processed/main/manifests/}.


% \subsubsection*{Stage 2}
% In our second stage, we ask the model to do some synthetic text generation. We need to train the model to rewrite the analysis into the FOMC Minutes. We have 1,392 Q\&A pairs in the dataset, covering the FOMC Minutes from Jan 2009 to Jan 2025. We randomly split it into 1,134 (80\%) training data, 139 (10\%) evaluation data, and 140 (10\%) test data. 




% For our second task, we 













